\section{A recap of standard logic and logical calculus}\label{sec:03-propositions-and-proofs:02-logic-recap:1}

\subsection{Recall}\label{sec:03-propositions-and-proofs:02-logic-recap:2}

Formal definitions cited in this section, gathered here for quick reference (full citations in the \hyperref[sec:root:bibliography:1]{Bibliography}):

\begin{itemize}
\tightlist
\item
  \textbf{Natural deduction.} After Gentzen (\cite{Gentzen1935}, ``Recherches sur la déduction logique,'' §II): each connective has an introduction rule (``une introduction (E)'') and an elimination rule (``une élimination (B)'') --- how to prove a formula built with that connective, and how to use one once it is available.
\item
  \textbf{Soundness.} ``\(\Gamma \vdash \varphi \Rightarrow \Gamma \models
  \varphi\)'' (\cite{VanDalen2013}, §2.5, Lemma 2.5.1). Brief: everything provable is true.
\item
  \textbf{Completeness.} ``\(\Gamma \vdash \varphi \Leftrightarrow \Gamma
  \models \varphi\)'' (\cite{VanDalen2013}, §2.5 Theorem 2.5.13 for propositional logic; §4.1 Theorem 4.1.3 for first-order logic; first-order completeness is originally due to Gödel, 1929/1930 --- not independently verified against Gödel's own paper, see \cite{VanDalen2013} for the identical statement independently verified here). Brief: everything true is provable.
\end{itemize}

The previous section introduced the Curry--Howard correspondence by translating logic directly into Lean types. It assumed that ``propositional logic,'' ``\(\vdash\),'' and ``natural deduction'' were at least half-familiar to the reader. If they were not, this section is the missing prerequisite --- skip ahead to \hyperref[sec:03-propositions-and-proofs:03-theorem-lemma:1]{§3} if propositional/first-order logic is already comfortable territory. It is a self-contained recap of standard mathematical logic --- the \emph{pre-Lean, pre-type-theory} version, exactly as it is presented in a first logic course --- so that §1's table has something concrete on its ``Logic'' side to refer back to. Nothing here mentions Lean, types, or programs; that translation was entirely §1's job, and \hyperref[sec:01-basics:05-pi-sigma-and-coc:1]{Chapter 1 §4} and \hyperref[sec:05-rigor-check:03-typing-rules-and-safety:1]{Chapter 5 §3} build the calculus those types compile down to. This section only fixes what the logic itself is.

\subsection{Propositional logic: syntax}\label{sec:03-propositions-and-proofs:02-logic-recap:3}

Fix a set of \textbf{propositional variables} (atomic statements) \(p, q, r,
\dots\). These stand for sentences whose internal structure is not analyzed --- ``it is raining,'' ``\(n\) is prime,'' anything with a definite truth value. \textbf{Formulas} are built from these using the \textbf{connectives}:

\[
\varphi ::= p \;\mid\; \top \;\mid\; \bot \;\mid\; \neg\varphi \;\mid\;
\varphi \wedge \varphi \;\mid\; \varphi \vee \varphi \;\mid\;
\varphi \Rightarrow \varphi
\]

read: a propositional variable, ``true,'' ``false,'' ``not \(\varphi\),'' ``\(\varphi\) and \(\varphi\),'' ``\(\varphi\) or \(\varphi\),'' ``\(\varphi\) implies \(\varphi\).'' This is pure syntax: a formula is just a string built by this grammar, nothing more. Whether a formula is \emph{true}, and whether it is \emph{provable}, are two separate questions, addressed next.

\subsection{Semantics: truth tables and validity}\label{sec:03-propositions-and-proofs:02-logic-recap:4}

A propositional variable's truth value is given by a \textbf{valuation} \(v : \{\text{variables}\} \to \{0, 1\}\) (an assignment of true or false to each atom). Every valuation extends uniquely to all formulas by the familiar truth tables:

{\small\begin{tabular}{
  >{\raggedright\arraybackslash}p{(\linewidth - 10\tabcolsep) * \real{0.1667}}
  >{\raggedright\arraybackslash}p{(\linewidth - 10\tabcolsep) * \real{0.1667}}
  >{\raggedright\arraybackslash}p{(\linewidth - 10\tabcolsep) * \real{0.1667}}
  >{\raggedright\arraybackslash}p{(\linewidth - 10\tabcolsep) * \real{0.1667}}
  >{\raggedright\arraybackslash}p{(\linewidth - 10\tabcolsep) * \real{0.1667}}
  >{\raggedright\arraybackslash}p{(\linewidth - 10\tabcolsep) * \real{0.1667}}
}
\toprule
\(\varphi\) & \(\psi\) & \(\varphi \wedge \psi\) & \(\varphi \vee \psi\) & \(\varphi \Rightarrow \psi\) & \(\neg \varphi\) \\
\midrule
0 & 0 & 0 & 0 & 1 & 1 \\
0 & 1 & 0 & 1 & 1 & 1 \\
1 & 0 & 0 & 1 & 0 & 0 \\
1 & 1 & 1 & 1 & 1 & 0 \\
\bottomrule
\end{tabular}}


A formula is a \textbf{tautology} (valid, written \(\models \varphi\)) if it comes out true under \emph{every} valuation. For example, \(p \vee \neg p\) (the \textbf{law of excluded middle}) and \(\neg\neg p \Rightarrow p\) (\textbf{double negation elimination}) are both tautologies: check every row of their truth tables and the result is always 1. This is the ``meaning-based'' side of logic. Truth is defined by checking every possible case, with no notion of proof or derivation involved at all.

\subsection{Proof theory: natural deduction}\label{sec:03-propositions-and-proofs:02-logic-recap:5}

\textbf{Provability}, by contrast, is a purely syntactic notion. A formula \(\varphi\) is \textbf{provable from hypotheses \(\Gamma\)} (a set of formulas), written \(\Gamma \vdash \varphi\), if there is a finite derivation of \(\varphi\) from \(\Gamma\) built out of a fixed, finite list of allowed \textbf{inference rules}. These are mechanical, symbol-pushing steps, checkable by an algorithm with no appeal to ``meaning'' at all. \textbf{Natural deduction} (Gentzen, 1934) is the standard system of such rules. Each connective gets an \textbf{introduction rule} (how to \emph{prove} a formula built with that connective) and an \textbf{elimination rule} (how to \emph{use} one once you have it). §1 already showed this pattern concretely (\texttt{⟨\_\allowbreak{}, \_\allowbreak{}⟩} introduces \texttt{∧}, \texttt{.left} eliminates it) without naming it. Writing \(\Gamma, \varphi\) for ``\(\Gamma\) together with the extra hypothesis \(\varphi\)'':

\[
\text{($\wedge$-intro)}\ \ \frac{\Gamma \vdash \varphi \qquad \Gamma \vdash \psi}
{\Gamma \vdash \varphi \wedge \psi}
\qquad\qquad
\text{($\wedge$-elim)}\ \ \frac{\Gamma \vdash \varphi \wedge \psi}{\Gamma \vdash \varphi}
\ \ \ \frac{\Gamma \vdash \varphi \wedge \psi}{\Gamma \vdash \psi}
\]

\[
\text{($\vee$-intro)}\ \ \frac{\Gamma \vdash \varphi}{\Gamma \vdash \varphi \vee \psi}
\ \ \ \frac{\Gamma \vdash \psi}{\Gamma \vdash \varphi \vee \psi}
\qquad\qquad
\text{($\vee$-elim)}\ \ \frac{\Gamma \vdash \varphi \vee \psi \qquad \Gamma, \varphi \vdash \chi \qquad \Gamma, \psi \vdash \chi}{\Gamma \vdash \chi}
\]

\[
\text{($\Rightarrow$-intro)}\ \ \frac{\Gamma, \varphi \vdash \psi}{\Gamma \vdash \varphi \Rightarrow \psi}
\qquad\qquad
\text{($\Rightarrow$-elim, modus ponens)}\ \ \frac{\Gamma \vdash \varphi \Rightarrow \psi \qquad \Gamma \vdash \varphi}{\Gamma \vdash \psi}
\]

\[
\text{($\neg$-intro)}\ \ \frac{\Gamma, \varphi \vdash \bot}{\Gamma \vdash \neg\varphi}
\qquad\qquad
\text{($\bot$-elim, ex falso)}\ \ \frac{\Gamma \vdash \bot}{\Gamma \vdash \varphi}
\]

\(\Rightarrow\)-intro reads: ``if, granting \(\varphi\) as an extra hypothesis, \(\psi\) can be derived, then (discharging that hypothesis) one may conclude \(\varphi \Rightarrow \psi\) outright.'' This is exactly the ordinary mathematical move ``assume \(\varphi\); \ldots{} ; therefore \(\varphi
\Rightarrow \psi\),'' turned into an explicit, checkable rule. It is \emph{exactly} what §1 identified with writing a Lean function \texttt{fun (hp : P) => ...}. Each rule above is stated once so it can be pointed to by name. The list need not be memorized --- the point is to recognize a ``natural deduction proof'' as a \emph{tree} built by chaining these rules, with leaves at hypotheses in \(\Gamma\) and its root at the conclusion \(\varphi\).

\textbf{Worked example: proving \(p \Rightarrow (q \Rightarrow p)\).} Assume \(p\) as a hypothesis, aiming to apply \(\Rightarrow\)-intro at the end. Within that, assume \(q\) too. The conclusion \(p\) is now already among the hypotheses, so it is derived for free. Discharge the \(q\)-hypothesis via \(\Rightarrow\)-intro to get \(q \Rightarrow p\), then discharge the \(p\)-hypothesis via \(\Rightarrow\)-intro again to get \(p \Rightarrow (q \Rightarrow p)\). As a derivation tree, with hypotheses listed to the left of \(\vdash\):

\[
\dfrac{\dfrac{p, q \vdash p}{p \vdash q \Rightarrow p}\ (\Rightarrow\text{-intro})}
{\vdash p \Rightarrow (q \Rightarrow p)}\ (\Rightarrow\text{-intro})
\]

\hyperref[sec:03-propositions-and-proofs:04-implication:1]{§3 (Implication)} names this exact formula ``implication is a function type'' and gives the corresponding Lean term directly: \texttt{fun hp => fun hq => hp}. The two are not just similar --- under Curry--Howard they are literally the same object, described twice.

\subsection{Soundness and completeness: proof theory meets semantics}\label{sec:03-propositions-and-proofs:02-logic-recap:6}

Two theorems connect the syntactic notion (\(\vdash\)) to the semantic one (\(\models\)). Together they justify treating ``provable'' and ``true in every case'' as interchangeable for propositional logic:

\begin{itemize}
\tightlist
\item
  \textbf{Soundness}: if \(\Gamma \vdash \varphi\) then \(\Gamma \models \varphi\). Everything the rules can derive really is true whenever the hypotheses are --- the rules never let you ``prove'' something false.
\item
  \textbf{Completeness} (Gödel, 1929, for first-order logic; the propositional case is elementary): if \(\Gamma \models \varphi\) then \(\Gamma \vdash
  \varphi\). Every semantically valid consequence \emph{does} have a natural deduction proof, so the rule list above, small as it is, is not missing anything.
\end{itemize}

Soundness is usually the easy direction to prove (check that each rule preserves truth); completeness is the harder theorem. Neither is used again in this book, but together they are the reason a working mathematician can trust that ``prove it'' and ``it is necessarily true'' describe the same territory for propositional (and first-order) logic. This guarantee stops holding once expressive enough systems are reached. (Gödel's \emph{incompleteness} theorems are a different and unrelated pair of results, despite the similar name --- they show arithmetic itself cannot be both complete and consistent.)

\subsection{First-order logic: adding quantifiers}\label{sec:03-propositions-and-proofs:02-logic-recap:7}

Propositional logic treats ``\(n\) is prime'' as one indivisible atom. \textbf{First-order logic} (also called predicate logic) opens that up: fix a domain of individuals, \textbf{predicates} \(P(x), Q(x, y), \dots\) ranging over it, and add two quantifiers to the grammar:

\[
\varphi ::= \cdots \;\mid\; \forall x.\, \varphi \;\mid\; \exists x.\, \varphi
\]

with natural deduction rules that generalize \(\wedge\)/\(\vee\)'s rules in the obvious way. \(\forall\)-intro requires proving \(\varphi\) for an \emph{arbitrary, otherwise-unconstrained} \(x\) (exactly ``let \(x\) be arbitrary; \ldots; therefore \(\forall x, \varphi\)'' from ordinary proof-writing). \(\exists\)-intro requires exhibiting one specific witness \(a\) and a proof of \(\varphi(a)\). And \(\exists\)-elim lets you reason from ``some \(x\) satisfies \(\varphi\)'' by naming an arbitrary such \(x\) and deriving the goal for it. This section's job so far has been to show that the \emph{quantifiers themselves}, and the rules governing them, are standard first-order logic with no Lean involved yet. The translation into Lean is deliberately left for the next paragraph, so that it is clear which half is ``ordinary logic already familiar to the reader'' and which half is ``Curry--Howard's doing.''

\textbf{First-order logic and Curry--Howard.} §1's table translated the \emph{propositional} connectives (\(\wedge, \vee, \Rightarrow, \neg\)) into type formers. Quantifiers extend the same table, and this is the one place where the translation genuinely needs \emph{dependent} types rather than ordinary ones, because \(P(x)\) --- the very thing being quantified over --- is a different proposition for each \(x\):

{\small\begin{tabular}{
  >{\raggedright\arraybackslash}p{(\linewidth - 4\tabcolsep) * \real{0.3333}}
  >{\raggedright\arraybackslash}p{(\linewidth - 4\tabcolsep) * \real{0.3333}}
  >{\raggedright\arraybackslash}p{(\linewidth - 4\tabcolsep) * \real{0.3333}}
}
\toprule
Logic & Type theory & Lean notation \\
\midrule
\(\forall x{:}\alpha,\ P(x)\) & dependent function (\(\Pi\)-)type & \texttt{∀ x, P x} \\
\(\exists x{:}\alpha,\ P(x)\) & dependent pair (\(\Sigma\)-)type & \texttt{∃ x, P x} \\
\(\forall\)-intro (arbitrary \(x\), derive \(\varphi\)) & \texttt{fun x => ...} & building a term of \texttt{∀ x, P x} \\
\(\forall\)-elim (instantiate at \(a\)) & function application & \texttt{(h : ∀ x, P x) a : P a} \\
\(\exists\)-intro (witness \(a\), proof \(h\)) & anonymous constructor & \texttt{⟨a, h⟩ : ∃ x, P x} \\
\(\exists\)-elim (unpack witness + proof) & pattern match / projection & \texttt{obtain ⟨a, h⟩ := ...}, \texttt{.1}/\texttt{.2} \\
\bottomrule
\end{tabular}}


Read the first row concretely: \(\forall x, P(x)\) becomes a \emph{dependent} function type precisely because its return type, \(P(x)\), depends on the very argument \(x\) being fed in. An ordinary (non-dependent) function type \texttt{α → β} would not be expressive enough, since \texttt{β} there is one fixed type, not one proposition per \texttt{x}. This is exactly \hyperref[sec:01-basics:03-dependent-types:1]{Chapter 1 §3}'s ``dependent types'' made concrete for the special case where the family being depended on happens to land in \texttt{Prop} instead of \texttt{Type}. \(\forall\)-elim is nothing more than ordinary function application: feed the function a specific \texttt{a}, get back a proof of \texttt{P a}. The natural-deduction rule and the programming-language operation are, again, not just similar but \emph{identical} --- the same fact \hyperref[sec:03-propositions-and-proofs:04-implication:1]{§3 (Implication)} already showed for modus ponens and plain function application.

With this table in hand, the whole of first-order natural deduction --- every rule stated in this section --- is visible as a special case of one simple idea: \emph{proofs are programs, and the specific shape of program a proof compiles to (pair, function, tagged choice, dependent function, dependent pair) is read off directly from the outermost connective or quantifier of the proposition being proved.} \hyperref[sec:01-basics:05-pi-sigma-and-coc:1]{Chapter 1 §4} makes this fully rigorous for the dependent case, by showing \(\Pi\) and \(\Sigma\) inside the calculus of constructions itself, rather than only stating the correspondence informally as this table does.

\subsection{Classical vs.~intuitionistic: the fork that matters for this book}\label{sec:03-propositions-and-proofs:02-logic-recap:8}

Every rule listed above is accepted by \textbf{both} classical and \textbf{intuitionistic} logic. The two differ on exactly one further principle, the \textbf{law of excluded middle}, \(\varphi \vee \neg\varphi\) (equivalently, double negation elimination \(\neg\neg\varphi \Rightarrow \varphi\)). Classical logic takes it as an extra axiom, valid for every \(\varphi\) regardless of whether you can exhibit a witness or decide the matter constructively. Intuitionistic logic --- the system natural deduction \emph{as given above} actually is, with no extra axiom added --- rejects it as a general principle: \(\varphi \vee \neg\varphi\) is not derivable from the rules above for an arbitrary \(\varphi\), only for specific \(\varphi\) that can actually be settled one way or the other. (\hyperref[sec:03-propositions-and-proofs:05-and-or-not:1]{§4 (And, Or, Not)}'s \href{https://lean-lang.org/doc/reference/latest/Tactic-Proofs/Tactic-Reference/}{\texttt{decide}} works because \texttt{1 = 2} happens to be \emph{decidable}, not because excluded middle is assumed.)

This is not just a side note: it is the precise reason Curry--Howard works as cleanly as it does. A type-theoretic proof term is a genuine, \emph{constructive} witness. A Lean proof of \(\exists x, P\, x\) computes to an actual pair \texttt{⟨a, h⟩} that could be inspected via \texttt{\#eval}, and that only makes sense for a logic where ``true'' means ``constructible'' --- which is exactly intuitionistic logic. (Lean's core logic is intuitionistic for precisely this reason. Mathlib freely adds classical excluded middle as an axiom for propositions where a witness is not needed, but the base calculus this book's Curry--Howard table describes in §1 does not include it.) Keep this fork in mind when reading \hyperref[sec:03-propositions-and-proofs:05-and-or-not:1]{§4 (And, Or, Not)}'s remark that Lean has ``no built-in law of excluded middle.'' It is this exact distinction, not an incidental implementation detail.

\subsection{References}\label{sec:03-propositions-and-proofs:02-logic-recap:9}

Full citations in the \hyperref[sec:root:bibliography:1]{Bibliography}. Formal definitions and verbatim quotes are gathered in Recall, above.

\begin{itemize}
\tightlist
\item
  Gentzen (\cite{Gentzen1935}), Section II ``Le calcul de la déduction naturelle'' --- natural deduction.
\item
  van Dalen (\cite{VanDalen2013}), §2.4 ``Natural Deduction'' (Definition 2.4.1, propositional rules), §2.5 ``Completeness'' (Lemma 2.5.1 Soundness, Theorem 2.5.13 Completeness), §3.8 ``Natural Deduction'' (first-order ∀-rules, Lemma 3.8.2), §4.1 ``The Completeness Theorem'' (Theorem 4.1.3) --- soundness, completeness.
\item
  Gödel --- \textbf{Gap:} the original 1930 completeness paper is not held in either notebook (see \texttt{NOTEBOOK-\allowbreak{}SOURCE-\allowbreak{}GAPS.md}); van Dalen's Theorem 4.1.3 above states the same result and is independently verified.
\item
  Pierce et al.~(\cite{PierceSF}) --- \textbf{Note:} only the \emph{Software Foundations} series homepage is available in the notebook, not chapter content, so the specific natural-deduction/classical-vs-intuitionistic treatment claimed here could not be verified verbatim.
\end{itemize}
